%!TEX root = ../main.tex
%!TEX program = xelatex
% Chapter 5: Words 1001 - 1250
\chapter{Nuanced Concepts}
\label{chap:chapter5}

{\small\sffamily\color{mutedtext} Abstract ideas, emotions, professional contexts, and multifaceted verbs.\par}

\vspace{0.4em}
\markboth{Chapter 5: Nuanced Concepts}{Words 1001--1250}

\begin{multicols}{2}
\vocentry{1001}{Операция}{ape'ratsija}{An operation; a surgery; a transaction}{Это секретная военная операция большой важности.}{It’s a secret military operation of great importance.}
\vocentry{1002}{Пространство}{prast'ranstvə}{Space, room}{Каждому члену семьи нужно личное пространство.}{Every member of the family needs their personal space.}
\vocentry{1003}{Спокойный}{spa'kojnyj}{Calm, serene}{У этой породы собак спокойный и ласковый характер.}{This dog breed has a calm and tender character.}
\vocentry{1004}{Одежда}{a'dezhda}{Clothes, clothing}{Одежда должна быть чистой и аккуратной.}{Clothes should be neat and clean.}
\vocentry{1005}{Кусок}{ku'sok}{A piece}{Для салата мне нужен большой кусок мягкого сыра.}{I need a big piece of soft cheese for the salad.}
\vocentry{1006}{Тема}{'tema}{A subject, a topic}{Тема этого рассказа мне очень близка.}{The subject of this short story is very close to my heart.}
\vocentry{1007}{Животный}{zhi'votnyj}{Animal (adj)}{Животный мир Сибири очень разнообразен.}{The animal world of Siberia is very diverse.}
\vocentry{1008}{Снимать}{sni'mat’}{To take off, to take down; to film}{Пожалуйста, не забывай снимать обувь, когда заходишь в дом.}{Please, don’t forget to take off your shoes when you enter the house.}
\vocentry{1009}{Искусство}{is'kustvə}{Art}{Настоящее искусство понятно всем и не нуждается в объяснениях.}{True art is clear to everyone and doesn’t need any explanations.}
\vocentry{1010}{Умный}{'umnyj}{Clever, smart, intelligent}{Он достаточно умный, чтобы не обращать внимания на эти сплетни.}{He’s clever enough not to pay attention to this gossip.}
\vocentry{1011}{Малый}{'malyj}{Small, little}{Этот закон поддерживает малый бизнес.}{This law supports small business.}
\vocentry{1012}{Курс}{kurs}{A course; a policy}{Я подписалась на этот курс и ни разу не пожалела об этом.}{I signed up for this course and I’ve never regretted it.}
\vocentry{1013}{Звонить}{zva'nit’}{To call (by phone); to ring}{Ты можешь звонить мне в любое время.}{You can call me any time.}
\vocentry{1014}{Очередной}{atchered'noj}{Next (in order), another}{Очередной этап соревнования состоится через неделю.}{The next step of the competition will take place in a week.}
\vocentry{1015}{Чудо}{'tchudə}{A miracle}{То, что мы выжили в той катастрофе, это настоящее чудо.}{The fact that we survived that disaster is a real miracle.}
\vocentry{1016}{Ощущение}{astchust'chenije}{A feeling, a sensation}{Мой дом дарит мне ощущение покоя и безопасности.}{My house gives me a feeling of peace and security.}
\vocentry{1017}{Крайне}{'krajne}{Extremely, very (much)}{Сотрудничество с этой компанией стало крайне негативным опытом для меня.}{The cooperation with that company was an extremely negative experience for me.}
\vocentry{1018}{Напоминать}{napami'nat’}{To remind; to resemble}{Я не могу всё время напоминать тебе о твоих обязанностях!}{I can’t remind you about your duties all the time!}
\vocentry{1019}{Придумать}{pri'dumat’}{To figure out; to make up}{Мы должны придумать, как подбодрить её.}{We must figure out how to cheer her up.}
\vocentry{1020}{Падать}{'padat’}{To fall; to decline}{Эксперты говорят, что цены будут продолжать падать.}{Experts say that the prices will go on falling.}
\vocentry{1021}{Здание}{'zdanije}{A building}{Они планируют снести это старое здание.}{They’re planning to pull down this old building.}
\vocentry{1022}{Свежий}{'svezhij}{Fresh (new or clean)}{Свежий кофе – это всё, что мне нужно, чтобы проснуться.}{Fresh coffee is all I need to wake up.}
\vocentry{1023}{Поговорить}{pagava'rit’}{To talk to, to have a conversation; to talk (to have a conversation that lasts a certain period of time; perfective)}{Могу я поговорить с вашим начальником?}{May I talk to your boss?}
\vocentry{1024}{Начинаться}{natchi'natsa}{To start, to begin}{Эссе должно начинаться со вступления.}{An essay should start with an introduction.}
\vocentry{1025}{Милиция}{mi'litsija}{Militia, police}{Милиция прибыла на место преступления первой.}{The militia was the first to arrive at the crime scene.}
\vocentry{1026}{Приказать}{prika'zat’}{To order, to command}{Я не могу приказать тебе остаться, но я был бы очень благодарен.}{I can’t order you to stay but I’d appreciate it a lot.}
\vocentry{1027}{Несмотря}{nesma'tr’a}{In spite of, despite (always used with preposition ‘на’)}{Несмотря на расстояние, они всегда поддерживают связь.}{In spite of the distance they always keep in touch.}
\vocentry{1028}{Оттуда}{at'tuda}{From there, out of there}{Давай поднимемся на крышу. Оттуда хорошо смотреть на звёзды.}{Let’s get up to the roof. It’s great to watch the stars from there.}
\vocentry{1029}{Засмеяться}{zasme'jatsa}{To laugh (means to start laughing)}{Я не мог не засмеяться, когда услышал твою шутку.}{I couldn’t help but laugh when I heard your joke.}
\vocentry{1030}{Волна}{val'na}{A wave, a tide}{Волна была такая огромная, что мы подумали, что это цунами.}{The wave was so huge that we thought it was a tsunami.}
\vocentry{1031}{Задача}{za'datcha}{A task, an assignment}{Эта задача требует комплексного подхода.}{This task demands a comprehensive approach.}
\vocentry{1032}{Касаться}{ka'satsa}{To concern; to touch}{Конфликты между родителями не должны касаться детей.}{The conflicts between parents mustn’t concern children.}
\vocentry{1033}{Старуха}{sta'ruha}{An old woman (colloquial)}{Люди говорили, что старуха умеет колдовать.}{People used to say that the old woman could do witchcraft.}
\vocentry{1034}{Войско}{'vojskə}{An army (military), troops}{Повстанцам удалось собрать большое войско.}{The insurgents managed to gather a great army.}
\vocentry{1035}{Срок}{srok}{A term; a deadline}{Председатель был избран на пятилетний срок.}{The chairman was elected for a five-year term.}
\vocentry{1036}{Ужас}{'uzhas}{Horror, terror}{Невозможно описать ужас, который она пережила.}{It’s impossible to describe the horror she experienced.}
\vocentry{1037}{Узкий}{'uzkij}{Narrow}{Мне не нравится эта квартира. Коридор слишком узкий и комнаты маленькие.}{I don’t like this apartment. The corridor is too narrow and the rooms are small.}
\vocentry{1038}{Материал}{mater’'jal}{Material, stuff}{Материал, из которого сделаны брюки, очень приятный на ощупь.}{The material the trousers are made of is very pleasant to touch.}
\vocentry{1039}{Приняться}{pri'n’atsa}{To get down to, to start doing}{У вас не так много времени, так что я бы советовал вам приняться за дело.}{You haven’t got much time, so I’d advise you to get down to business.}
\vocentry{1040}{Нормальный}{nar'mal’nyj}{Normal; sane}{В этой ситуации он повел себя как любой нормальный подросток.}{In that situation he behaved like any normal teenager.}
\vocentry{1041}{Крик}{krik}{A cry, a scream}{Её громкий крик был слышен даже на втором этаже.}{Her loud cry was heard even on the second floor.}
\vocentry{1042}{Знание}{'znanije}{Knowledge}{Её глубокое знание и понимание предмета помогли ей получить повышение.}{Her deep knowledge and understanding of the subject helped her get a promotion.}
\vocentry{1043}{Трое}{'troje}{Three (in a group together)}{Все трое детей простудились после прогулки под дождём.}{All the three kids caught a cold after the walk in the rain.}
\vocentry{1044}{Четвёртый}{tchet'vjortyj}{Fourth}{Это его четвёртый гол за сезон.}{It’s his fourth goal in the season.}
\vocentry{1045}{Шум}{shum}{A noise}{Шум в квартире соседей мешал мне уснуть.}{The noise in the neighbors’ apartment prevented me from falling asleep.}
\vocentry{1046}{Заставить}{zas'tavit’}{To make, to force}{Никто не сможет заставить меня изменить моё мнение.}{No one will be able to make me change my mind.}
\vocentry{1047}{Автомат}{afto'mat}{A machine gun; a vending machine}{Я понятия не имею, как работает автомат и чем он отличается от винтовки.}{I’ve got no idea how a machine gun works and how it’s different from a rifle.}
\vocentry{1048}{Резко}{'reskə}{Sharply, abruptly, harshly}{Она ответила резко и грубо.}{She answered sharply and rudely.}
\vocentry{1049}{Тётя}{'t’ot’a}{An aunt; a woman (colloquial)}{Моя тётя часто участвует в волонтерских проектах.}{My aunt often takes part in volunteer projects.}
\vocentry{1050}{Впечатление}{fpetchat'lenije}{An impression}{Невозможно произвести первое впечатление дважды.}{It’s impossible to make the first impression twice.}
\vocentry{1051}{Пауза}{'pauza}{A pause}{В нашем разговоре возникла долгая неловкая пауза.}{There was a long and awkward pause in our conversation.}
\vocentry{1052}{Глубина}{glubi'na}{Depth}{Глубина этой реки не достаточно большая для кораблей.}{The depth of this river is not big enough for ships.}
\vocentry{1053}{Доллар}{'dolar}{A dollar}{Доллар – одна из самых популярных мировых валют.}{The dollar is one of the most popular world currencies.}
\vocentry{1054}{Сотня}{'sotn’a}{A hundred}{У неё была сотня причин расстаться с ним.}{She had a hundred reasons to break up with him.}
\vocentry{1055}{Виноватый}{vina'vatyj}{Guilty}{У собаки был такой виноватый вид, что я не мог злиться на неё.}{The dog had such a guilty look that I couldn’t be angry with it.}
\vocentry{1056}{Студент}{stu'dent}{A student}{Я студент, поэтому я не могу работать на полную ставку.}{I am a student, which is why I can’t work full-time.}
\vocentry{1057}{Боевой}{baje'voj}{Battle, combat (adj)}{Дети играли в индейцев и придумали специальный боевой клич.}{The children were playing Cowboys and Indians and made up a special battle cry.}
\vocentry{1058}{Очко}{atch'ko}{A point (in a game); blackjack}{Игрок получает одно очко за каждый правильный ответ.}{A player gets one point for every correct answer.}
\vocentry{1059}{Нож}{nosh}{A knife}{Этот нож слишком тупой, чтобы резать им мясо.}{This knife is too blunt to cut meat with.}
\vocentry{1060}{Подниматься}{padni'matsa}{To rise, to get up}{Уровень воды в реке продолжал подниматься, и ситуация становилась опасной.}{The level of water in the river was continuing to rise and the situation was getting dangerous.}
\vocentry{1061}{Мировой}{mira'voj}{World, global, worldwide}{Возможно ли было предвидеть мировой финансовый кризис?}{Was it possible to foresee the world financial crisis?}
\vocentry{1062}{Предмет}{pred'met}{An object; a subject, a topic}{При раскопках учёные обнаружили странный предмет и сейчас изучают его.}{During the excavations the scientists found a strange object and are studying it now.}
\vocentry{1063}{Плохой}{pla'hoj}{Bad}{Это плохой пример. Он ничего не объясняет.}{It’s a bad example. It doesn’t explain anything.}
\vocentry{1064}{Майор}{ma'jor}{A major (military)}{Старый майор выглядел серьёзным и недружелюбным.}{The old major looked serious and unfriendly.}
\vocentry{1065}{Ударить}{u'darit’}{To hit, to strike, to punch}{Если ты не перестанешь, мне придётся тебя ударить.}{If you don’t stop, I’ll have to hit you.}
\vocentry{1066}{Разуметься}{razu'metsa}{Very rarely used in the initial form. Most often used as part of a set expression ‘Само собой разумеется’ - ‘It goes without saying’.}{Само собой разумеется, что полное выздоровление займёт много времени.}{It goes without saying that full recovery will take much time.}
\vocentry{1067}{Точный}{'tochnyj}{Exact, precise}{Нам нужен точный анализ всех параметров.}{We need an exact analysis of all the parameters.}
\vocentry{1068}{Всякий}{'vsjakij}{Any, every; all sorts of}{Всякий раз, когда я слышу эту песню, я вспоминаю детство.}{Anytime I hear this song I recall my childhood.}
\vocentry{1069}{Деятельность}{'dejatel’nəst’}{Activity}{Его общественная деятельность приносит много пользы другим людям.}{His social activity does a lot of good to other people.}
\vocentry{1070}{Инженер}{inzhe'ner}{An engineer}{Мой дядя – опытный инженер, и его часто приглашают протестировать новое оборудование.}{My uncle is an experienced engineer and he’s often invited to test new equipment.}
\vocentry{1071}{Программа}{prag'rama}{A program/programme}{Программа обучения состоит из двух частей.}{The training program consists of two parts.}
\vocentry{1072}{Костюм}{kast’'um}{A suit; a costume}{Этот костюм не подходит мне по размеру.}{This suit doesn’t fit me.}
\vocentry{1073}{Женский}{'zhenskij}{Female, feminine, woman’s, women’s}{Контральто – это тип классического женского певческого голоса.}{A contralto is a type of classical female singing voice.}
\vocentry{1074}{Крепкий}{'krepkij}{Tough, strong; firm, solid}{Твой брат крепкий парень. Эта работа для него.}{Your brother is a tough guy. This job is for him.}
\vocentry{1075}{Театр}{te'atr}{A theater}{Театр – мой любимый вид искусства.}{Theater is my favorite kind of art.}
\vocentry{1076}{Встречаться}{vstre'chatsa}{To date; to see each other}{Они начали встречаться 3 месяца назад и уже собираются пожениться.}{They started dating three months ago and they are going to marry already.}
\vocentry{1077}{Различный}{raz'lichnyj}{Different; various, diverse}{Мы не можем поместить всех детей в одну группу. У них у всех различный уровень подготовленности.}{We can’t place all the children in one group. They’re all at different levels of training.}
\vocentry{1078}{Хвост}{hvost}{A tail}{У моей кошки длинный пушистый хвост.}{My cat has a long fluffy tail.}
\vocentry{1079}{Танк}{tank}{A tank}{Сегодня этот боевой танк – музейный экспонат.}{Today this battle tank is a museum exhibit.}
\vocentry{1080}{Социальный}{satsi'al’nyj}{Social}{Этот социальный центр поддерживает иммигрантов.}{This social center supports immigrants.}
\vocentry{1081}{Мёртвый}{'m’ortvyj}{Dead; a deadman}{Моя дочь ужасно боится насекомых. Даже мёртвый жук может напугать её.}{My daughter is terribly scared of insects. Even a dead beetle can frighten her.}
\vocentry{1082}{Попытаться}{papy'tatsa}{To attempt, to try}{Я хочу хотя бы попытаться помириться с ней.}{I want at least to attempt to make up with her.}
\vocentry{1083}{Явно}{'javnə}{Obviously, clearly}{Она явно что-то скрывает от нас.}{She’s obviously hiding something from us.}
\vocentry{1084}{Журнал}{zhur'nal}{A magazine}{Журнал мод – это единственная книга на её полке.}{A fashion magazine is the only book on her shelf.}
\vocentry{1085}{Конь}{kon’}{A horse (a male horse)}{Мне кажется, тот чёрный конь будет быстрее остальных.}{It seems to me that the black horse will be faster than the rest.}
\vocentry{1086}{Схватить}{shva'tit’}{To grasp, to seize}{Я успел схватить мальчика за руку прежде, чем он упал с лестницы.}{I managed to grasp the boy by the hand before he fell off the ladder.}
\vocentry{1087}{Вдоль}{vdol’}{Along}{Мы гуляли вдоль реки и разговаривали о прошлом.}{We were walking along the river and talking about the past.}
\vocentry{1088}{Стихи}{sti'hi}{Poems}{Когда я был подростком, я, как и многие другие, пробовал писать стихи.}{When I was a teenager I, like many others, tried to write poems.}
\vocentry{1089}{Степень}{'stepen’}{An extent, a degree}{Степень участия в проекте может быть разной. Главное, чтобы все внесли свой вклад.}{The extent of participation in the project may be different. The main thing is that everyone adds value.}
\vocentry{1090}{Отказаться}{atka'zatsa}{To refuse; to give up}{Я не в настроении ехать на экскурсию. Могу я отказаться?}{I’m not in the mood to go on excursion. Can I refuse?}
\vocentry{1091}{Вытащить}{'vytastchit’}{To pull out, to drag out}{Мне пришлось схватить щенка за уши, чтобы вытащить его из воды.}{I had to grasp the puppy by the ears to pull it out of water.}
\vocentry{1092}{Предлагать}{predla'gat’}{To offer; to suggest}{Наша компания планирует предлагать услуги маркетинга.}{Our company is planning to offer marketing services.}
\vocentry{1093}{Фигура}{fi'gura}{A figure}{У тебя идеальная фигура. Что ты делаешь, чтобы держать себя в форме?}{You’ve got a perfect figure. What do you do to keep in shape?}
\vocentry{1094}{Ключ}{kl’utch}{A key}{Я потерял ключ от входной двери, и мне пришлось заменить замок.}{I’ve lost the key to the front door and had to replace the lock.}
\vocentry{1095}{Чёрт}{tchort}{Devil (very often used in set expressions like ‘Чёрт с тобой!’ - “To hell with you!”)}{Он такой хитрый, что, похоже, даже сам чёрт не сможет его обмануть.}{He’s so cunning that, apparently, even the devil himself won’t be able to deceive him.}
\vocentry{1096}{Отправиться}{at'pravitsa}{To go on, to go to; to set off (perfective)}{Я хотел отправиться в круиз, но мне пришлось всё отменить из-за болезни отца.}{I wanted to go on a cruise but I had to cancel everything because of my father’s illness.}
\vocentry{1097}{Побежать}{pabe'zhat’}{To start running}{Она хотела побежать за автобусом, но поняла, что всё равно не догонит его.}{She wanted to start running after the bus but realized she would not catch up with it anyway.}
\vocentry{1098}{Вагон}{va'gon}{A car (on a train)}{Спальный вагон был душным и неудобным.}{The sleeping car was stuffy and uncomfortable.}
\vocentry{1099}{Культура}{kul’tura}{Culture}{Культура каждой нации по-своему уникальна.}{The culture of every nation is unique in its own way.}
\vocentry{1100}{Курить}{ku'rit’}{To smoke}{Ты должен бросить курить, иначе тебя не примут в команду.}{You must give up smoking, otherwise they won’t accept you in the team.}
\vocentry{1101}{Боец}{ba'jets}{A fighter; a soldier}{Это маленькая девочка настоящий боец. У неё было три операции!}{This little girl is a real fighter. She has had three surgeries!}
\vocentry{1102}{Западный}{'zapadnyj}{Western}{Западный мир никогда не поймёт восточный и наоборот.}{The Western world will never understand the Eastern one and vice versa.}
\vocentry{1103}{Фильм}{fil’m}{A film, a movie}{Этот исторический фильм вдохновил меня на написание новой статьи.}{This historical film inspired me to write a new article.}
\vocentry{1104}{Восемь}{'vosem’}{Eight}{В этом классе всего восемь учеников.}{There are only eight pupils in this class.}
\vocentry{1105}{Открывать}{atkry'vat’}{To open}{Я прошу тебя не открывать подарки до твоего дня рождения.}{I ask you not to open the presents before your birthday.}
\vocentry{1106}{Зайти}{zaj'ti}{To drop by (perfective)}{Я хочу поговорить с тобой. Можно мне зайти после работы?}{I want to talk to you. May I drop by after work?}
\vocentry{1107}{Немедленно}{ne'medlennə}{Immediately}{Вы должны немедленно отвезти ребёнка в больницу!}{You must immediately take the child to the hospital!}
\vocentry{1108}{Замечать}{zame'tchat’}{To notice, to note}{Теперь я начинаю замечать, что ты сильно изменился.}{Now I’m beginning to notice that you’ve changed a lot.}
\vocentry{1109}{Пыль}{pyl’}{Dust}{Нам придётся долго убирать этот старый дом. Пыль повсюду!}{We’ll have to clean this old house for a long time. Dust is everywhere!}
\vocentry{1110}{Закончить}{za'kontchit’}{To finish, to end (perfective)}{Мне нужно ещё немного времени, чтобы закончить книгу.}{I need some more time to finish the book.}
\vocentry{1111}{Мягкий}{'mjahkij}{Soft; mild}{Какой мягкий пирог! Он просто тает во рту!}{What a soft cake! It’s just melting in the mouth!}
\vocentry{1112}{Труба}{tru'ba}{A pipe, a tube; a trumpet}{Я только что заметил, что эта труба протекает. Мы должны вызвать сантехника.}{I’ve just noticed that this pipe is leaking. We must call a plumber.}
\vocentry{1113}{Редкий}{'redkij}{Rare}{Это редкий вид бабочек. Он занесён в Красную книгу.}{It’s a rare species of butterfly. It’s listed in the Red Book.}
\vocentry{1114}{Миллион}{mili'on}{A million}{Что бы ты делал, если бы выиграл миллион долларов?}{What would you do if you won a million dollars?}
\vocentry{1115}{Следить}{sled’'it’}{To follow, to watch; to spy}{Мы будем следить за событиями на соревнованиях.}{We are going to follow the events at the competitions.}
\vocentry{1116}{Масса}{'masa}{Mass (weight); mass (a lot of)}{Какова точная масса этой машины? Наше оборудование сможет поднять её?}{What’s the exact weight of this car? Will our equipment be able to lift it?}
\vocentry{1117}{Пятьдесят}{p’at’de's’at}{Fifty}{Ты читала “Пятьдесят оттенков серого”?}{Have you read “The Fifty Shades of Grey”?}
\vocentry{1118}{Обстоятельство}{absta'jatel’stvə}{A circumstance}{Это обстоятельство помешало нам закончить проект раньше.}{That circumstance prevented us from finishing the project earlier.}
\vocentry{1119}{Решиться}{re'shitsa}{To dare}{Она не могла решиться рассказать нам правду, боясь, что мы не поймём её.}{She couldn’t dare to tell us the truth because she was afraid we wouldn’t understand her.}
\vocentry{1120}{Одновременно}{adnavre'mennə}{Simultaneously}{Я слышал несколько голосов одновременно.}{I could hear a few voices simultaneously.}
\vocentry{1121}{Броситься}{'brositsa}{To rush, to dash}{Мне хотелось броситься к ней и обнять её.}{I wanted to rush to her and embrace her.}
\vocentry{1122}{Пошлый}{'poshlyj}{Vulgar}{Его пошлый анекдот смутил девушку.}{His vulgar joke embarrassed the girl.}
\vocentry{1123}{Круглый}{'kruglyj}{Round (adj)}{Посреди комнаты стоит маленький круглый стол.}{There’s a small round table standing in the middle of the room.}
\vocentry{1124}{Штаб}{shtab}{A headquarters (military)}{Штаб фронта был расположен на севере страны.}{The headquarters of the front was situated in the North of the country.}
\vocentry{1125}{Слух}{sluh}{Hearing; a rumor}{Слух может ухудшаться с возрастом.}{Hearing can get worse with age.}
\vocentry{1126}{Голый}{'golyj}{Bare, naked, nude}{В комнате не было даже ковра, только голый пол.}{There wasn’t a carpet in the room, only bare floor.}
\vocentry{1127}{Закрыть}{zak'ryt’}{To close, to shut (perfective)}{Можешь закрыть окно? Холодно.}{Can you close the window? It’s cold.}
\vocentry{1128}{Линия}{'linija}{A line}{Красная линия разделила страницу на две части.}{A red line divided the page into two parts.}
\vocentry{1129}{Лодка}{'lodka}{A boat}{Это лодка моего дедушки. Он на ней рыбачит.}{It’s my granddad’s boat. He goes fishing on it.}
\vocentry{1130}{Наоборот}{naaba'rot}{Vice versa, the other way round}{Это ты должен извиняться передо мной, а не наоборот.}{It’s you who should apologize to me and not vice versa.}
\vocentry{1131}{Даль}{dal’}{Distance, far}{Она задумчиво смотрела в недостижимую даль.}{She was looking into the distance thoughtfully.}
\vocentry{1132}{Председатель}{predse'datel’}{A chairman}{Председатель комитета предложил начать голосование.}{The chairman of the committee suggested to start the vote.}
\vocentry{1133}{Понравиться}{pan'ravitsa}{To like (the person who likes takes the dative case) (perfective)}{Эта книга должна тебе понравиться. Это новый детектив.}{You must like this book. It’s a new detective story.}
\vocentry{1134}{Обратиться}{abra'titsa}{To turn to, to address (perfective)}{Тебе нужно обратиться к специалисту. Это слишком сложный вопрос.}{You should turn to a specialist. This matter is too complicated.}
\vocentry{1135}{Период}{pe'riəd}{A period; an epoch}{Это был сложный период их отношений, но они смогли сохранить их.}{That was a hard period of their relationship but they managed to save it.}
\vocentry{1136}{Организм}{arga'nism}{An organism, a body}{Человеческий организм – это удивительный механизм.}{The human organism is an amazing mechanism.}
\vocentry{1137}{Отдел}{at'del}{A department}{Отдел продаж находится в другом здании.}{The sales department is in a different building.}
\vocentry{1138}{Весна}{ves'na}{Spring}{Ранняя весна обычно холодная и дождливая.}{Early spring is usually cold and rainy.}
\vocentry{1139}{Обед}{a'bed}{Dinner, lunch}{Что у нас на обед? Я умираю от голода!}{What have we got for dinner? I’m starving!}
\vocentry{1140}{Значение}{zna'tchenije}{A meaning}{У этого слова есть ещё одно значение.}{This word has one more meaning.}
\vocentry{1141}{Дно}{dno}{Bottom}{Дно океана – это загадочное место.}{The bottom of the ocean is a mysterious place.}
\vocentry{1142}{Фотография}{fota'grafija}{A photo}{Эта семейная фотография мне очень дорога.}{This family photo is very dear to me.}
\vocentry{1143}{Некий}{'nekij}{Some, a certain}{Некий господин Иванов хочет видеть Вас.}{Some Mr. Ivanov wants to see you.}
\vocentry{1144}{Успех}{us'peh}{Success}{Успех проекта зависит от качества сотрудничества.}{The success of the project depends on the quality of collaboration.}
\vocentry{1145}{Грязный}{'grjaznyj}{Dirty, muddy}{Почему ковёр такой грязный? Мне придётся отдать его в химчистку.}{Why is the carpet so dirty? I’ll have to take it to the dry-cleaner’s.}
\vocentry{1146}{Естественный}{jes'testvennyj}{Natural}{Это твой естественный цвет волос?}{Is it your natural hair color?}
\vocentry{1147}{Велеть}{ve'let’}{To order, to command (literary, archaic)}{Если хочешь, я могу велеть ей вернуть твои вещи.}{If you want, I can order her to give your things back.}
\vocentry{1148}{Сцена}{'stsena}{A stage; a scene}{Сцена – это её жизнь. Она отличная актриса.}{The stage is her life. She’s an excellent actress.}
\vocentry{1149}{Билет}{bi'let}{A ticket}{У меня есть свободный билет на матч. Пойдёшь со мной?}{I have a spare ticket to the match. Are you going with me?}
\vocentry{1150}{Состав}{sas'tav}{Content, composition}{Я всегда обращаю внимание на состав продуктов питания.}{I always pay attention to the content of foods.}
\vocentry{1151}{Знаменитый}{zname'nityj}{Famous}{Знаменитый пианист даёт сегодня концерт в нашем городе.}{Today a famous pianist is giving a concert in our city.}
\vocentry{1152}{Лишний}{'lishnyj}{Extra; superfluous, excess}{Я знала, что ты придёшь, и подготовила лишний стул.}{I knew you’d come and prepared an extra chair.}
\vocentry{1153}{Честь}{tchest’}{Honor}{Эта улица названа в честь героя войны.}{This street is named in honor of a war hero.}
\vocentry{1154}{Иметься}{i'metsa}{There be, there is/are, to be, to be available}{В любом городе должна иметься хотя бы одна больница.}{There should be at least one hospital in any town.}
\vocentry{1155}{Смех}{smeh}{Laughter, laugh}{В парке развлечений детский смех был слышен повсюду.}{In the amusement park children’s laughter was heard everywhere.}
\vocentry{1156}{СССР}{ɛs ɛs ɛs 'ɛr}{USSR}{СССР был антирелигиозным государством.}{USSR was an anti-religious state.}
\vocentry{1157}{Деревянный}{dere'vjannyj}{Wooden}{Наш дачный домик деревянный, поэтому летом в нём прохладно.}{Our country house is wooden, which is why in summer it’s cool inside.}
\vocentry{1158}{Опустить}{apus'tit’}{To put down; to omit}{Полицейский приказал грабителю опустить пистолет.}{The policeman ordered the robber to put down the gun.}
\vocentry{1159}{Родиться}{ra'ditsa}{To be born}{Я бы хотел родиться в другую эпоху.}{I wish I were born in a different epoch.}
\vocentry{1160}{Способ}{'sposəb}{A means, a way}{Интернет – самый популярный способ общения среди молодежи.}{The Internet is the most popular means of communication among the youth.}
\vocentry{1161}{Спасти}{spas'ti}{To save, to rescue}{Им удалось спасти альпиниста только через три дня.}{They managed to save the climber in only three days.}
\vocentry{1162}{Зрение}{'zrenije}{Eyesight, vision}{У неё очень плохое зрение. Она ничего не видит без очков.}{She’s got very poor eyesight. She doesn’t see anything without glasses.}
\vocentry{1163}{Английский}{an'glijskij}{English}{Английский язык самый популярный из всех мировых языков.}{The English language is the most popular among the world languages.}
\vocentry{1164}{Отдельный}{at'del’nyj}{Separate, detached; individual}{У этой комнаты есть отдельный вход.}{This room has a separate entrance.}
\vocentry{1165}{Извинить}{izvi'nit’}{To excuse, to pardon, to forgive (for)}{Вам придётся извинить меня, но мне пора идти. Спасибо за отличный вечер!}{You’ll have to excuse me but it’s time for me to go. Thank you for the great night!}
\vocentry{1166}{Громкий}{'gromkij}{Loud}{Меня разбудил громкий раскат грозы.}{I was woken up by a loud peal of thunder.}
\vocentry{1167}{Объявить}{abja'vit’}{To announce, to declare}{Мы решили объявить о помолвке.}{We’ve decided to announce the engagement.}
\vocentry{1168}{Запад}{'zapad}{West}{Запад страны более развит экономически.}{The west of the country is more developed economically.}
\vocentry{1169}{Германия}{ger'manija}{Germany}{Германия уделяет большое внимание проблемам окружающей среды.}{Germany pays great attention to the problems of the environment.}
\vocentry{1170}{Обращаться}{abra'stchatsa}{To turn to, to address}{Я не люблю обращаться к врачам за помощью и редко хожу в больницу.}{I don’t like to turn to doctors for help and rarely go to hospital.}
\vocentry{1171}{Население}{nase'lenije}{Population}{Население страны растёт недостаточно быстро.}{The population of the country doesn’t grow fast enough.}
\vocentry{1172}{Беда}{be'da}{A trouble, a misfortune}{Та страшная беда навсегда изменила его жизнь.}{That terrible trouble changed his life forever.}
\vocentry{1173}{Спешить}{spe'shit’}{To hurry, to be in a hurry}{Нам нужно спешить. Такси уже ждёт.}{We need to hurry. The taxi is waiting already.}
\vocentry{1174}{Еда}{je'da}{Food; a meal}{Здоровая еда – это ключ к долгой жизни.}{Healthy food is the key to a long life.}
\vocentry{1175}{Поздно}{'poznə}{Late}{Уже поздно. Детям пора ложиться спать.}{It’s late already. It’s time for the kids to go to bed.}
\vocentry{1176}{Специальный}{spetsi'alnyj}{Special, custom}{Я выиграл специальный приз конкурса.}{I won the special prize of the contest.}
\vocentry{1177}{Наблюдать}{nabl’u'dat’}{To watch, to observe}{Я люблю наблюдать за птицами. Они удивительные создания.}{I like to watch birds. They’re amazing creatures.}
\vocentry{1178}{Еврей}{ev'rej}{A Jew, a Hebrew}{Этот пожилой еврей – очень хороший ювелир.}{This elderly Jew is a very good goldsmith.}
\vocentry{1179}{Ой}{oj}{Oh, oops, auch}{Ой, я наступила тебе на ногу. Я не нарочно.}{Oh, I’ve stepped on your foot. That wasn’t on purpose.}
\vocentry{1180}{Суметь}{su'met’}{To be able, to manage, to succeed}{Ты должен суметь убедить её поехать с нами.}{You must be able to persuade her to go with us.}
\vocentry{1181}{Обернуться}{aber'nutsa}{To turn around; to result in}{Не успел я обернуться, как собака схватила мой бутерброд и выбежала из комнаты.}{Hardly had I managed to turn around when the dog seized my sandwich and ran out of the room.}
\vocentry{1182}{Министр}{mi'nistr}{A minister}{Министр иностранных дел прибыл с официальным визитом.}{The Minister of Foreign Affairs arrived on an official visit.}
\vocentry{1183}{Фирма}{'firma}{A firm, a company}{Их фирма обанкротилась во время кризиса.}{Their firm went bankrupt during the crisis.}
\vocentry{1184}{Проснуться}{pras'nutsa}{To wake up (perfective)}{Я не могу проснуться без чашки крепкого кофе.}{I can’t wake up without a cup of strong coffee.}
\vocentry{1185}{Организация}{argani'zatsija}{An organisation}{ООН – это влиятельная международная организация.}{The UN is an influential international organization.}
\vocentry{1186}{Принцип}{'printsip}{A principle}{Я не могу понять принцип работы этого механизма.}{I can’t understand the principle of work of this mechanism.}
\vocentry{1187}{Направление}{naprav'lenije}{A direction}{Я думаю, мы выбрали правильное направление развития.}{I think we’ve chosen the right direction of development.}
\vocentry{1188}{Гораздо}{ga'razdə}{Much, far, by far}{Мне кажется, что книга гораздо интереснее фильма.}{It seems to me that the book is much more interesting than the movie.}
\vocentry{1189}{Мастер}{'master}{A master, an expert}{Этот ремесленник – настоящий мастер. Его работы уникальны.}{This craftsman is a real master. His works are unique.}
\vocentry{1190}{Подождать}{padazh'dat’}{To wait (perfective)}{Ваши фотографии ещё не готовы. Вам придётся подождать.}{Your photos are not ready yet. You’ll have to wait.}
\vocentry{1191}{Светлый}{'svetlyj}{Light (shade)}{Я выбираю этот светлый плащ. Тёмный мне не подходит.}{I choose this light raincoat. The dark one doesn’t suit me.}
\vocentry{1192}{Ошибка}{a'shibka}{A mistake}{Эта ошибка испортила весь результат.}{This mistake ruined the whole result.}
\vocentry{1193}{Высота}{vysa'ta}{Height}{Высота этого небоскрёба просто поражает.}{The height of this skyscraper is just stunning.}
\vocentry{1194}{Существо}{sustchest'vo}{A creature}{Синий кит – самое большое существо на планете.}{The blue whale is the biggest creature on the planet.}
\vocentry{1195}{Заговорить}{zagava'rit’}{To begin to talk}{Он не мог осмелиться даже заговорить с ней.}{He couldn’t dare even to begin talking with her.}
\vocentry{1196}{Мальчишка}{mal’'tchishka}{A boy, a kid, a brat}{Ты просто избалованный мальчишка!}{You’re just a spoiled boy!}
\vocentry{1197}{Пользоваться}{'pol’zavatsa}{To use; to enjoy (success for example)}{Ты можешь пользоваться любыми вещами, которые тебе нужны. Чувствуй себя как дома.}{You can use anything you need. Make yourself at home.}
\vocentry{1198}{Научный}{'nautchnyj}{Scientific}{Давай попробуем применить научный подход.}{Let’s try to apply a scientific approach.}
\vocentry{1199}{Повторять}{pavta'rjat’}{To repeat; to revise}{При изучении иностранного языка очень полезно повторять каждое слово вслух.}{It’s very useful to repeat each word aloud when learning a foreign language.}
\vocentry{1200}{Яркий}{'jarkij}{Bright}{Мне нравится яркий узор на твоей майке.}{I like the bright pattern on your T-shirt.}
\vocentry{1201}{Несколько}{'neskal’ka}{A few, some}{У меня есть несколько вопросов. Вы можете уделить мне немного времени?}{I’ve got a few questions. Can you devote some time to me?}
\vocentry{1202}{Учить}{u'tchit’}{To learn, to study; to teach}{Я не могу учить новые слова, не записывая их.}{I can’t learn new words without writing them down.}
\vocentry{1203}{Взяться}{'vzjatsa}{To undertake}{Я не смогу взяться за новый проект, пока не закончу этот.}{I won’t be able to undertake a new project until I finish this one.}
\vocentry{1204}{Бегать}{'begat’}{To run}{Пожалуйста, перестань бегать по дому и успокойся.}{Please, stop running about the house and calm down.}
\vocentry{1205}{Господь}{gas'pod’}{Lord, God}{Господь милостив! Он поможет!}{Lord is merciful! He’ll help!}
\vocentry{1206}{Общественный}{abst'chestvennyj}{Public, common; social}{Я трачу довольно много денег на общественный транспорт.}{I spend quite a lot of money on public transport.}
\vocentry{1207}{Устроить}{ust'roit’}{To arrange, to organize}{Хотите комнату с видом на море? Я могу это устроить.}{Do you want a room with a sea view? I can arrange that.}
\vocentry{1208}{Дышать}{dy'shat’}{To breathe}{Хочу жить там, где смогу дышать чистым воздухом.}{I want to live where I can breathe clean air.}
\vocentry{1209}{Лечь}{letch’}{To lie down}{Тебе лучше лечь и оставаться в постели некоторое время.}{You’d better lie down and stay in bed for some time.}
\vocentry{1210}{США}{sɛ shɛ 'a}{The USA}{США расположены между Мексикой и Канадой.}{The USA is situated between Mexico and Canada.}
\vocentry{1211}{Камера}{'kamera}{A camera (photography); a chamber; a cell (in prison)}{Мне нужна новая камера, чтобы делать фотографии более высокого качества.}{I need a new camera to take photos of a better quality.}
\vocentry{1212}{Победа}{pa'beda}{Victory}{Победа в чемпионате сделала его всемирно знаменитым.}{The victory in the championship made him world famous.}
\vocentry{1213}{Учёный}{u'tchjonyj}{Learned, academic}{Мой учёный друг лучше энциклопедии, но он такой рассеянный!}{My learned friend is better than an encyclopedia but he’s so absent-minded!}
\vocentry{1214}{Состоять}{sasta'jat’}{To consist (of); to be a member of}{Экзамен будет состоять из трёх частей, включая практическое задание.}{The exam will consist of three parts including the practical task.}
\vocentry{1215}{Заходить}{zaha'dit’}{To enter, to come in (to); to drop by}{Я не разрешаю заходить в мою комнату без стука.}{I don’t allow anyone to enter my room without knocking.}
\vocentry{1216}{Человечество}{tchela'vetchestvə}{Mankind, humankind, humanity}{Сегодня человечество утратило многие традиционные ценности.}{Today mankind has lost many traditional values.}
\vocentry{1217}{Фраза}{'fraza}{A phrase}{Его фраза прозвучала очень грубо.}{His phrase sounded very rude.}
\vocentry{1218}{Внимательно}{vni'matel’nə}{Attentively}{Постарайтесь слушать внимательно. Я не буду повторять дважды.}{Try to listen attentively. I won’t repeat twice.}
\vocentry{1219}{Считаться}{stchi'tatsa}{To be considered (impersonal)}{Как такое поведение может считаться нормальным?}{How can such behavior be considered normal?}
\vocentry{1220}{Вероятно}{vera'jatnə}{Probably, very likely}{Вы, вероятно, ещё не слышали последние новости.}{You probably haven’t heard the latest news yet.}
\vocentry{1221}{Соседний}{sa'sednij}{Neighboring, next to}{Соседний дом продаётся. Ты не знаешь, сколько они хотят за него?}{The neighboring house is for sale. Do you know how much they want for it?}
\vocentry{1222}{Замок}{za'mok}{A lock}{На этом сейфе кодовый замок. Его очень сложно открыть.}{There’s a coded lock on this safe. It’s very hard to open.}
\vocentry{1223}{Рыжий}{'ryzhyj}{Red, ginger (of hair)}{Мой толстый рыжий кот такой ленивый! Он даже не знает, как ловить мышей.}{My fat ginger cat is so lazy! He doesn’t even know how to catch mice.}
\vocentry{1224}{Ездить}{'jezdit’}{To go (by horse or vehicle), to ride, to drive}{Он любит ездить в экзотические страны. Он уже посетил больше десятка.}{He likes going to exotic countries. He’s already visited more than a dozen.}
\vocentry{1225}{Господи}{'gospadi}{Lord, God (as a vocative)}{О, Господи! Что же мы теперь будем делать?}{Oh, Lord! What are we going to do now?}
\vocentry{1226}{Остановить}{astana'vit’}{To stop (perfective)}{Нужно делать всё возможное, чтобы остановить глобальное потепление.}{We should do everything possible to stop the global warming.}
\vocentry{1227}{Встретиться}{'vstretitsa}{To meet (with); to come across}{Я хочу встретиться с ней и обсудить детали контракта.}{I want to meet with her and discuss the details of the contract.}
\vocentry{1228}{Явиться}{ja'vitsa}{To show up, to turn up}{Послушай, ты не можешь явиться на собеседование в шортах.}{Listen, you can’t show up to a job interview in shorts.}
\vocentry{1229}{Рынок}{'rynək}{A marketplace; a market}{Давай сходим на рынок в субботу. Я хочу купить свежей рыбы к празднику.}{Let’s go to the marketplace on Saturday. I want to buy some fresh fish for the holiday.}
\vocentry{1230}{Тяжело}{tjazhe'lo}{Hard, difficultly; heavily}{Тяжело совмещать семью и карьеру.}{It’s hard to combine a family and a career.}
\vocentry{1231}{Злой}{zloj}{Evil, mean}{Злой волшебник жил глубоко в лесу.}{The evil magician lived deep in the woods.}
\vocentry{1232}{Клуб}{klub}{A club}{В нашем районе открылся новый ночной клуб.}{They’ve opened a new night club in our district.}
\vocentry{1233}{Привезти}{prives'ti}{To bring (by vehicle)}{Эти строительные материалы слишком тяжёлые. Мы не сможем привезти их на своей машине.}{These building materials are too heavy. We won’t be able to bring them in our car.}
\vocentry{1234}{Платить}{pla'tit’}{To pay}{Жители города отказались платить новый налог.}{The citizens refused to pay the new tax.}
\vocentry{1235}{Кость}{kost’}{A bone}{Рентген показал, что кость сломана.}{The X-ray showed that the bone was broken.}
\vocentry{1236}{Дикий}{'dikij}{Wild}{Этот волк – дикий зверь. Я не думаю, что ты сможешь приручить его.}{This wolf is a wild beast. I don’t think you’ll be able to tame it.}
\vocentry{1237}{Личность}{'lichnəst’}{Personality; a person, an individual}{Он действительно уникальная личность. Я раньше никогда не встречал таких людей.}{He’s got a really unique personality. I’ve never met such people before.}
\vocentry{1238}{Столица}{sta'litsa}{A capital}{Любая столица – это центр культурной и политической жизни страны.}{Any capital is the center of cultural and political life of the country.}
\vocentry{1239}{Спасть}{spast’}{To break, to fall (to decrease)}{Прогноз погоды говорит, что жара должна спасть к концу месяца.}{The weather forecast says that the heat should break by the end of the month.}
\vocentry{1240}{Здоровье}{zda'rovje}{Health}{Моё здоровье заметно улучшилось по сравнению с прошлым годом.}{My health has improved notably in comparison with the previous year.}
\vocentry{1241}{Здравствовать}{'zdrastvəvat’}{To prosper, to thrive}{Мы надеемся, что вы будете здравствовать ещё долгие годы!}{We hope you’ll prosper for many years to come!}
\vocentry{1242}{Попытка}{pa'pytka}{A try, an attempt}{Хорошая попытка! Но ты опять не угадал.}{A nice try! But you haven’t guessed it again.}
\vocentry{1243}{Позволять}{pazva'ljat’}{To allow, to permit}{Как ты можешь позволять своим детям так с тобой разговаривать?}{How can you allow your children to talk to you like that?}
\vocentry{1244}{Наиболее}{nai'boleje}{The most, most of all}{Это наиболее эффективный метод лечения.}{It’s the most efficient method of treatment.}
\vocentry{1245}{Постепенно}{paste'pennə}{Gradually, step by step}{Изменения будут внедряться постепенно.}{The changes are going to be implemented gradually.}
\vocentry{1246}{Пятый}{'pjatyj}{Fifth}{Май – это пятый месяц года.}{May is the fifth month of the year.}
\vocentry{1247}{Определённый}{aprede'ljonnyj}{Certain, definite}{В определённый момент мы просто перестали общаться.}{At a certain moment we just stopped communicating.}
\vocentry{1248}{Читатель}{tchi'tatel’}{A reader}{Дорогой читатель, эта книга полна приключений.}{Dear reader, this book is full of adventures.}
\vocentry{1249}{Шутка}{'shutka}{A joke}{Шутка была такая смешная, что мы катались от смеха.}{The joke was so funny that we were rolling with laughter.}
\vocentry{1250}{Исторический}{ista'ritcheskij}{Historical; historic}{Исторический центр города привлекает много туристов.}{The historical center of the city attracts lots of tourists.}
\end{multicols}
